%% =====================================================================
%% curved_dunn_raw_direct_sum_platonic.tex
%%
%% Raw direct-sum curved-Dunn comparison for class M.
%%
%% Doctrine. Genus 1 is governed by the separate twisted tensor
%% criterion. For g >= 2, curved-Dunn H^2-vanishing is a conditional
%% transport theorem: modular-bootstrap H^2-acyclicity plus the
%% degree-2 comparison map imply curved-Dunn acyclicity. KZB boundary
%% regularity supplies analytic clutching only under boundary formal
%% type and mixed-channel corner data; it is not the proof of
%% curved-Dunn H^2-vanishing.
%%
%% Location: chapters/theory/.
%% Inputs: chapters/theory/curved_dunn_higher_genus.tex
%%         chapters/theory/modular_swiss_cheese_operad.tex
%%         standalone/curved_dunn_two_complex_bridge.tex
%%         standalone/stokes_gap_kzb_regularity.tex
%% =====================================================================

\chapter{Raw Direct-Sum Curved Dunn at Class M}
\label{ch:curved-dunn-raw-direct-sum-platonic}

\section*{Genus and ambient}

The curved-Dunn comparison has three regimes. Genus \(0\) is the
uncurved Dunn tensor product. Genus \(1\) is the elliptic twisted
tensor criterion of
\autoref{prop:genus1-twisted-tensor-product}: a twisting cochain
\(\tau_{1}=\mathrm{Mon}(R)\) must satisfy the curved tensor-product
Maurer--Cartan equation in the selected coderived ambient. Genus
\(g \ge 2\) is the conditional transport theorem of
\autoref{thm:curved-dunn-H2-vanishing-all-genera}: modular-bootstrap
\(H^{2}\)-acyclicity and the degree-\(2\) comparison map
\(\Phi\) are both required.

The raw direct sum is the bounded algebraic object obtained before
passing to a pro-object, a \(J\)-adic completion, or the
weight-completed category. It is the correct place to see the
finite-weight obstruction. It is not a replacement for the
conditional completed theorem.

\section{Raw complex and harmonic discrepancy}
\label{sec:setup-harmonic-discrepancy}

Fix a chirally Koszul algebra \(A\) on a smooth projective curve
\(X\) at non-critical level, with modular characteristic
\(\kappa(A)\). Let \(\omega_{g}\) denote the Arakelov form on
\(\overline{\cM}_{g,n}\), normalized as in
\autoref{ch:curved-dunn-higher-genus}.

\begin{definition}[Curved-Dunn twisting-cochain complex, raw form]
\label{def:curved-dunn-complex-raw}
For \(g \ge 2\), the \emph{raw direct-sum curved-Dunn
twisting-cochain object} is
\[
 \mathrm{TCo}^{\bullet}_{g,\oplus}(A)
 \;=\;
 \bigoplus_{r \ge 0}
 \mathrm{Hom}\!\Bigl(
 \mathbf{B}\bigl(E_{1}^{\mathrm{tr}}\bigr)_{[r]},\;
 \mathrm{gr}^{g}\mathbf{B}^{\mathrm{mod}}_{r}\bigl(
 \mathrm{SC}^{\mathrm{ch,top}}\bigr)\Bigr)
\]
in \(\mathrm{Ch}(\mathrm{Vect})\). The convolution differential is
\(d_{0}\). The curvature insertion
\[
 m_{g}=\kappa(A)\,\omega_{g}\wedge(-)
\]
is filtration-compatible, but the raw operator
\(D_{g,\oplus}=d_{0}+m_{g}\) is not assumed to be a differential.
\end{definition}

\begin{definition}[Weight filtration and harmonic discrepancy]
\label{def:harmonic-discrepancy}
The descending weight filtration on the raw object is
\[
F^{r}\mathrm{TCo}^{\bullet}_{g,\oplus}
 =
 \bigoplus_{s \ge r}
 \mathrm{Hom}\!\bigl(
 \mathbf{B}(E_{1}^{\mathrm{tr}})_{[s]},\,
 \mathrm{gr}^{g}\mathbf{B}^{\mathrm{mod}}_{s}\bigr).
\]
The \emph{raw harmonic discrepancy at weight \(r\)} is the
associated-graded component
\[
 \delta_{r}^{\mathrm{harm}}(A,g)
 :=
 D_{g,\oplus}^{2}\Big|_{F^{r}/F^{r+1}}
 \in
 \mathrm{Hom}\!\bigl(
 \mathbf{B}(E_{1}^{\mathrm{tr}})_{[r]},\,
 \mathrm{gr}^{g}\mathbf{B}^{\mathrm{mod}}_{r+2}\bigr).
\]
If \(D_{g,\oplus}^{2}=0\), every
\(\delta_{r}^{\mathrm{harm}}\) vanishes. A single non-zero
\(\delta_{r}^{\mathrm{harm}}\) rules out a raw direct-sum
vanishing theorem.
\end{definition}

\begin{lemma}[Low-weight harmonic discrepancy for class \(\mathrm{M}\)]
\label{lem:low-weight-discrepancy-class-M}
On the Virasoro class-\(\mathrm{M}\) line, the quartic contact
invariant is
\[
 S_{4}(\mathrm{Vir}_{c})=\frac{10}{c(5c+22)}.
\]
For the raw genus-\(2\) Virasoro object,
\[
 \delta_{4}^{\mathrm{harm}}(\mathrm{Vir}_{c},2)
 =
 S_{4}(\mathrm{Vir}_{c})\,\omega_{2}^{\wedge 2}
 =
 \frac{10}{c(5c+22)}\,\omega_{2}^{\wedge 2}.
\]
At generic \(c\), this class is non-zero whenever
\(\omega_{2}^{\wedge2}\) survives in the chosen cohomological
model.
\end{lemma}

\begin{proof}
The local landscape records
\(S_{4}(\mathrm{Vir}_{c})=10/[c(5c+22)]\) on the Virasoro line
(\autoref{sec:rosetta-stone}; see also the worked class-\(\mathrm{M}\)
examples). The raw associated-graded square of \(D_{2,\oplus}\) sees
the first quartic contact term at weight \(4\). Thus the displayed
component is the product of the Virasoro quartic invariant with the
genus-\(2\) curvature square. At \(c\notin\{0,-22/5\}\) the scalar is
non-zero; the remaining condition is exactly non-vanishing of the
geometric class in the cohomological model under consideration.
\end{proof}

\section{Raw direct-sum obstruction}
\label{sec:raw-direct-sum-failure}

\begin{theorem}[Raw direct-sum curved Dunn at class \(\mathrm{M}\)]
\label{thm:curved-dunn-class-M-raw-direct-sum-fails}
The bounded raw direct-sum ambient does not carry a hypothesis-free
curved-Dunn \(H^{2}\)-vanishing theorem. Already for
\(A=\mathrm{Vir}_{c}\) and \(g=2\),
\[
 D_{2,\oplus}^{2}\Big|_{F^{4}/F^{5}}
 =
 \frac{10}{c(5c+22)}\,\omega_{2}^{\wedge 2},
\]
so the raw square is non-zero at generic \(c\) whenever
\(\omega_{2}^{\wedge2}\neq0\).
\end{theorem}

\begin{proof}
This is the weight-four component of
\autoref{lem:low-weight-discrepancy-class-M}. A raw theorem would
force \(D_{2,\oplus}^{2}=0\), hence
\(\delta_{4}^{\mathrm{harm}}(\mathrm{Vir}_{c},2)=0\). The displayed
generic Virasoro value contradicts that conclusion.
\end{proof}

\begin{remark}[Direct-sum obstruction and completed ambients]
\label{rem:direct-sum-obstruction-is-ambient-choice}
The raw obstruction is not a proof of completed vanishing. It is a
diagnostic: the direct-sum object has retained a finite-weight
quartic contact term which must be killed by an actual comparison
primitive in the completed theory. The primitive is supplied only
under the hypotheses of
\autoref{thm:curved-dunn-H2-vanishing-all-genera}.
\end{remark}

\section{Completed comparison}
\label{sec:pro-ambient-reconstitution}

\begin{definition}[Pro-ambient curved-Dunn object]
\label{def:pro-curved-dunn}
The pro-ambient raw object at genus \(g\) is the inverse system
\[
 \mathrm{pro}\text{-}\mathrm{TCo}^{\bullet}_{g}(A)
 =
 \bigl\{\mathrm{TCo}^{\bullet}_{g,\le N}(A)\bigr\}_{N\ge0},
 \qquad
 \mathrm{TCo}^{\bullet}_{g,\le N}(A)
 =
 \bigoplus_{r\le N}
 \mathrm{Hom}\!\bigl(
 \mathbf{B}(E_{1}^{\mathrm{tr}})_{[r]},\,
 \mathrm{gr}^{g}\mathbf{B}^{\mathrm{mod}}_{r}\bigr),
\]
with transition maps the weight-truncation quotients
\(\pi_{N+1,N}\). The differential before curvature correction is
the induced convolution differential \(d_{0}\). A curved
strictification at level \(N\) is additional data, not a formal
consequence of truncation.
\end{definition}

\begin{lemma}[Finite truncation comparison criterion]
\label{lem:finite-truncation-is-complex}
Fix \(g\ge2\) and \(N\ge0\). Assume the modular-bootstrap complex
has \(H^{2}=0\) in genus \(g\) through weight \(N\), and assume the
degree-\(2\) comparison map
\[
 \Phi_{\le N}\colon
 (\fg^{A}_{\mathrm{mod},\le N},d_{\Theta^{(0)}})
 \longrightarrow
 (\mathrm{TCo}^{\bullet}_{g,\le N}(A),d_{0})
\]
identifies the modular-bootstrap and curved-Dunn obstruction
classes and induces an isomorphism on \(H^{2}\). Then every
degree-\(2\) curved-Dunn obstruction on
\(\mathrm{TCo}^{\bullet}_{g,\le N}(A)\) is \(d_{0}\)-exact.
Choosing a primitive gives a strictified truncated curved operator.
Without these two hypotheses, no finite-truncation strictification is
asserted.
\end{lemma}

\begin{proof}
The proof is the low-degree comparison argument of
\autoref{prop:modular-bootstrap-to-curved-dunn-bridge} restricted to
weights \(\le N\). A degree-\(2\) curved-Dunn cocycle lifts
along \(H^{2}(\Phi_{\le N})\) to a modular-bootstrap class. The
modular-bootstrap \(H^{2}\)-acyclicity makes that lift exact, and
applying \(\Phi_{\le N}\) gives a \(d_{0}\)-primitive on the
curved-Dunn side. Adding the chosen primitive is precisely the
finite-level strictification datum.
\end{proof}

\begin{theorem}[Conditional pro-ambient curved-Dunn \(H^{2}\)-vanishing]
\label{thm:curved-dunn-pro-ambient-H2-zero}
Let \(A\) be chirally Koszul at non-critical level and let
\(g\ge2\). Assume:
\begin{enumerate}
\item the modular-bootstrap \(H^{2}\)-acyclicity criterion holds in
genus \(g\);
\item the degree-\(2\) comparison criterion of
\autoref{prop:modular-bootstrap-to-curved-dunn-bridge} holds in
genus \(g\), compatibly with the weight truncations;
\item the inverse system of degree-\(1\) cohomology groups satisfies
the Mittag-Leffler condition.
\end{enumerate}
Then
\[
 H^{2}\bigl(\mathrm{pro}\text{-}\mathrm{TCo}^{\bullet}_{g}(A)\bigr)
 =0.
\]
\end{theorem}

\begin{proof}
By \autoref{lem:finite-truncation-is-complex}, each finite
truncation has vanishing degree-\(2\) obstruction under the stated
comparison and modular-bootstrap hypotheses. The transition maps are
weight quotients, and the Mittag-Leffler hypothesis kills the
\(\lim^{1}\) term in the Milnor exact sequence
\[
0\to
\lim\nolimits^{1}H^{1}(\mathrm{TCo}^{\bullet}_{g,\le N})
\to
H^{2}(\mathrm{pro}\text{-}\mathrm{TCo}^{\bullet}_{g})
\to
\lim H^{2}(\mathrm{TCo}^{\bullet}_{g,\le N})
\to0.
\]
The left term vanishes by Mittag-Leffler and the right term vanishes
by the finite-truncation comparison; hence the middle term is zero.
\end{proof}

\section{\(J\)-adic topological ambient}
\label{sec:j-adic-ambient}

\begin{definition}[Positive-weight ideal \(J\) and \(J\)-adic object]
\label{def:j-adic-complex}
Let
\[
J=\bigoplus_{r\ge1}\mathrm{TCo}^{\bullet,[r]}_{g}(A)
\subset \mathrm{TCo}^{\bullet}_{g,\oplus}(A)
\]
be the positive-weight ideal. The \(J\)-adic topological object is
\[
 \mathrm{TCo}^{\bullet}_{g,J\text{-adic}}(A)
 =
 \varprojlim_{N}
 \mathrm{TCo}^{\bullet}_{g,\oplus}(A)/J^{N+1}.
\]
It identifies with the same weight-truncation inverse system as
\(\mathrm{pro}\text{-}\mathrm{TCo}^{\bullet}_{g}(A)\) when the
finite-weight pieces and transition maps agree.
\end{definition}

\begin{theorem}[Conditional \(J\)-adic curved-Dunn \(H^{2}\)-vanishing]
\label{thm:curved-dunn-J-adic-H2-zero}
Under the hypotheses of
\autoref{thm:curved-dunn-pro-ambient-H2-zero}, and assuming the
\(J\)-adic quotients identify with the weight truncations, one has
\[
 H^{2}\bigl(\mathrm{TCo}^{\bullet}_{g,J\text{-adic}}(A)\bigr)=0.
\]
\end{theorem}

\begin{proof}
The \(J^{N+1}\)-quotient is the weight-\(\le N\) truncation. Thus
the \(J\)-adic object and the pro-object have the same inverse
system in degrees relevant to \(H^{2}\). Apply
\autoref{thm:curved-dunn-pro-ambient-H2-zero}.
\end{proof}

\section{Comparison of completed ambients}
\label{sec:three-ambients-agree}

\begin{proposition}[Completed ambient comparison in degree \(2\)]
\label{prop:three-ambients-equivalent}
Assume the finite-weight condition, the \(J\)-adic/pro
identification, and the degree-\(2\) comparison hypotheses of
\autoref{thm:curved-dunn-pro-ambient-H2-zero}. Then the pro-object,
the \(J\)-adic object, and the weight-completed curved-Dunn object
have the same degree-\(2\) cohomology. This common \(H^{2}\) vanishes
under the same hypotheses. No hypothesis-free full quasi-isomorphism
of completed ambients is asserted.
\end{proposition}

\begin{proof}
The pro and \(J\)-adic objects are the same inverse system under the
stated finite-weight identification. The map to the
weight-completed object is the completion map used in
\autoref{ch:curved-dunn-higher-genus}. In degree \(2\), its
acyclicity statement is exactly the conditional transport theorem:
modular-bootstrap \(H^{2}\)-acyclicity plus the degree-\(2\)
comparison. Thus the three degree-\(2\) groups agree in the
conditional range and vanish there.
\end{proof}

\section{Boundary KZB input}
\label{sec:kzb-boundary-input-raw}

\begin{remark}[KZB regularity and Stokes data]
KZB boundary composition is separate from curved-Dunn
\(H^{2}\)-vanishing. At a simple node with plumbing coordinate
\(\hbar\), the polar part is logarithmic:
\[
 \nabla^{\mathrm{KZB}}
 =
 d-\frac{A_{1}}{\hbar}\,d\hbar+\Omega_{\mathrm{reg}},
 \qquad
 A_{1}=\frac{\Omega_{\mathrm{node}}}{k+h^{\vee}}.
\]
The singularity is regular singular and has Newton slope \(0\).
It is not a Stokes source.

Stokes data enter only on boundary strata with non-zero irregular
slope after a meromorphic formal type
\(\Lambda(t)=\sum_{m=1}^{d}A_{m}t^{-m}\) is supplied, possibly after
finite ramification. The formal type determines Stokes directions
and Stokes groups; the Stokes matrices are additional analytic
monodromy data. At two-edge corners, edgewise formal types do not
imply coherence. One must impose the mixed-channel zero-curvature
condition appearing in
\autoref{prop:kzb-boundary-formal-type-edge-corner-reduction} and
\autoref{thm:irregular-singular-kzb-regularity}.
\end{remark}

\begin{corollary}[Direct-sum obstruction and completed replacement]
\label{cor:platonic-reconstitution-of-item-20}
The raw Virasoro weight-four witness rules out a hypothesis-free
bounded-direct-sum theorem. The completed pro/\(J\)-adic/
weight-completed conclusion is the conditional one: for \(g\ge2\)
it requires modular-bootstrap \(H^{2}\)-acyclicity and the
degree-\(2\) comparison map; for \(g=1\) it requires the separate
twisted tensor criterion. Boundary KZB regularity supplies analytic
composition only under compatible formal type and mixed-channel
corner hypotheses.
\end{corollary}

\begin{proof}
The first sentence is \autoref{thm:curved-dunn-class-M-raw-direct-sum-fails}.
The completed \(g\ge2\) statement is
\autoref{thm:curved-dunn-pro-ambient-H2-zero},
\autoref{thm:curved-dunn-J-adic-H2-zero}, and
\autoref{prop:three-ambients-equivalent}. The genus-\(1\) statement
is \autoref{prop:genus1-twisted-tensor-product}. The boundary KZB
claim is the preceding remark, which records the hypotheses of
\autoref{thm:irregular-singular-kzb-regularity}.
\end{proof}
